\section{Involution between the residue and quotient coordinates}
\label{sec:dualidad-visible-memoria}

The Euclidean division used in \(\mathcal A_9^\#\) writes, for an integer in
the nonadic chart,

\[
n=9Q+r,
\qquad r\in\{1,\ldots,9\}.
\]

The coordinate \(r\) is the observable representative, and \(Q\) records the
number of boundary crossings. To formulate an involution interchanging
the two coordinates, one must pass to an extension stable under the
interchange. Denote by \(\bar r\in\{0,\ldots,8\}\) the residual class,
with \(\bar r=0\) when the representative is \(r=9\), and consider

\[
\mathscr D_{\mathrm{dual}}=\mathbb Z^2\times\mathbb R_{>0}.
\]

The two integers of \(\mathscr D_{\mathrm{dual}}\) are not restricted to the canonical section of
representatives; this enlargement is necessary because a quotient \(Q\) may
fall outside the interval \(0,\ldots,8\) when occupying the first coordinate.

\begin{definition}[Radial energy form]
For \(R_0>0\) and \((m,w)\in\mathbb Z^2\), define
\[
E_{R_0}(m,w)=\frac{m^2}{R_0^2}+w^2R_0^2.
\]
The duality transformation is
\[
\mathsf T(m,w,R_0)=(w,m,R_0^{-1}).
\]
\end{definition}

\begin{theorem}[Residue--Quotient Duality]
\label{thm:dualidad-visible-memoria}
The map \(\mathsf T\) is an involution of \(\mathscr D_{\mathrm{dual}}\) and preserves
the energy form in the sense
\[
E_{R_0}(m,w)=E_{R_0^{-1}}(w,m).
\]
\end{theorem}

\begin{proof}
The second map of \(\mathsf T\) exchanges of new \(m,w\) e
invierte by second time \(R_0\), by lo that
\[
\mathsf T^2(m,w,R_0)=(m,w,R_0).
\]
Moreover,
\[
E_{R_0^{-1}}(w,m)
=\frac{w^2}{R_0^{-2}}+m^2R_0^{-2}
=w^2R_0^2+\frac{m^2}{R_0^2}
=E_{R_0}(m,w).\qedhere
\]
\end{proof}

The insertion of a nonadic state into the extended domain is

\[
(r,Q)\longmapsto(\bar r,Q,R_0).
\]

The transform \((Q,\bar r,R_0^{-1})\) need not belong to the
section \(0\leq\bar r\leq8\). For this reason, \(\mathsf T\) is an
exact involution of the enlarged domain, not an endomorphism of the residual atlas
of 81 positions. The fixed point of the parameter transformation is
\(R_0=1\); there the exchange of the two coordinates preserves the same
quadratic form.

The identity admits an operatorial realization. Let
\(\mathcal H_{\rm lat}=\ell^2(\mathbb Z^2)\), with orthonormal basis
\(\{|m,w\rangle\}\). Define the diagonal operator

\[
H(R_0)|m,w\rangle=E_{R_0}(m,w)|m,w\rangle
\]

in its natural domain, and the unitary operator

\[
U_T|m,w\rangle=|w,m\rangle.
\]

\begin{corollary}[Unitary Equivalence of Reciprocal Radii]
\label{cor:dualidad-unitaria}
For every \(R_0>0\),
\[
U_TH(R_0)U_T^{-1}=H(R_0^{-1}).
\]
\end{corollary}

\begin{proof}
On a basis vector,
\[
U_TH(R_0)U_T^{-1}|m,w\rangle
=E_{R_0}(w,m)|m,w\rangle.
\]
The \cref{thm:dualidad-visible-memoria} gives
\(E_{R_0}(w,m)=E_{R_0^{-1}}(m,w)\), which is the action of
\(H(R_0^{-1})\).
\end{proof}

This conjugation has the mathematical form of radius duality in a circular compactification, in units in which the string scale has been normalized to one \cite{Giveon1994}. The further physical identification assigns \(m\) to the momentum number, \(w\) to the winding number, and \(R_0\) to the compactification radius. The theorem does not introduce oscillators, fields, supersymmetry, or a string tension; it proves the lattice involution and the unitary equivalence of the declared diagonal operators.
